\documentclass[12pt]{article}
\usepackage[utf8]{inputenc}
\usepackage[T1]{fontenc}
\usepackage{hyperref}
\title{Puoc Language Resources for K-12 Curriculum Development}
\author{Compiled by Claude Sonnet 5.5 (claude-sonnet-5-5) for Andr\'as Kornai}
\date{}
\begin{document}
\maketitle

\section{Introduction}
Puoc (ISO 639-3 \texttt{puo}, Glottocode puoc1238) is an Austroasiatic language of the Khmuic branch, spoken in northwestern Vietnam (Son La and Lai Chau provinces) and Laos. It is also known as Ksingmul or Xinh Mun. Sources found give about 27,000 speakers and describe it as vigorous. Web search found only general reference pages; no school materials, primers, or native-language publications were found. This dossier therefore lists reference material only.

\section{Primary Educational Resources}
No primers, textbooks, or school programmes for Puoc were found in search results.

\section{Linguistic and Reference Materials}
\begin{itemize}
\item Ksingmul language (Wikipedia overview, names, varieties): \url{https://en.wikipedia.org/wiki/Ksingmul_language}
\item Ksingmul people (Wikipedia, ethnographic overview): \url{https://en.wikipedia.org/wiki/Ksingmul_people}
\item Khmuic languages (classification context): \url{https://ipfs.aleph.im/ipfs/QmXoypizjW3WknFiJnKLwHCnL72vedxjQkDDP1mXWo6uco/wiki/Khmuic.html}
\item An ANU Open Research repository item appeared in results (apparently a Khmuic bibliography); its content was not inspected: \url{https://openresearch-repository.anu.edu.au/server/api/core/bitstreams/a0cc836c-0a78-4a2b-b8e3-ac119c2f4bae/content}
\end{itemize}

\section{Religious and Community Resources}
None found.

\section{Digital and Audio Resources}
Wikitongues has a language page for Puoc; whether it holds recordings was not verified: \url{https://wikitongues.org/languages/puo/}

\section{Curriculum Development Context}
Little documented material exists. Any curriculum would need to begin with community consultation and new documentation (orthography, word lists, recordings), in the context of Vietnamese-medium schooling.

\end{document}